% MAS1004 Assignment 2: long report.
%
% You fill this in with your agent, in your repository. AGENTS.md tells the
% agent how. Red text in the PDF is what still has to be replaced.
%
% - Pictures go in figures/. Copy them from results/, which is not in git.
% - Outputs go in outputs/: copy experiments.csv, LOG.md and the outputs of
%   check_submission.py and curate.py count here. \outputfile{outputs/...}
%   prints a file exactly as it is. Every number in this report has to come
%   from such an output.
%
% To make the PDF: compress this folder into a zip, and on Overleaf choose
% New Project, Upload Project. Set Menu, Compiler to XeLaTeX, and Recompile.

\documentclass[11pt]{article}
\usepackage[ko]{mas1004}

\reporttitle{Assignment 2: Long report}
\studentname{김소연}
\studentid{20261179}

\begin{document}
\makeheader

Leaderboard nickname: \texttt{slowwbluee}. Track A 마지막 release: \texttt{track-a-3},
``every city drawn equally often''. Track B: \fillin{release 없음}.

% The sections below are a suggestion. Change, merge or add sections as you
% like; README.md, "The long report", says what a reader should be able to see.

% ---------------------------------------------------------------------------
\section*{Track A}

\begin{center}\footnotesize
\begin{tabular}{lp{0.44\linewidth}rrr}
\toprule
tag & 이전 실행에서 바뀐 것 & seed 0 & seed 1 & 평균 \\
\midrule
baseline & 없음. 시작 레시피, 학습 사진 44,257장 전부 & 0.5028 & 0.4879 & 0.4954 \\
street & \texttt{sample\_weights}: 출처별로 다른 횟수 (commons $\times$1, youtube $\times$2, osv5m $\times$4) & 0.4857 & 0.4812 & 0.4835 \\
drop\_street & 레시피는 시작 레시피로 되돌리고, 학습 사진 6,964장을 뺌 (37,293장 남음) & 0.4841 & 0.4918 & 0.4880 \\
balanced & \texttt{sample\_weights}: 다섯 도시를 똑같이 뽑음 (사진은 44,257장 전부) & 0.5038 & 0.5123 & 0.5081 \\
\bottomrule
\end{tabular}

{\footnotesize 표 1. 레시피마다의 검증 점수. 평균은 두 시드 행의 평균이다.}
\end{center}

모든 실행은 같은 조건에서 쟀다. 점수는 도시별 정답률 다섯 개의 평균이고, 학습 사진
수는 249,984장으로 고정이며(예산은 레시피가 정하지 않는다), 한 번에 한 가지만 바꾸고,
바꾼 것마다 시드 0과 시드 1로 두 번 학습해 두 값이 같은 방향으로 움직이는지 보았다.

\subsection*{baseline: 시작점과 잡음의 크기}

시작 레시피(batch 64, RandomResizedCrop + 좌우반전 + ColorJitter, AdamW
$3\times10^{-4}$, OneCycle, label smoothing 0.1)를 학습 사진 44,257장 전부로 두 번
학습했다. 레시피는 그대로 두고 시드만 0과 1로 바꾼 이유는 시작점을 잡고, 시드만
바꿨을 때 점수가 얼마나 흔들리는지 알고 싶었기 때문이다.

두 값은 0.5028과 0.4879로 0.0149 차이가 났다. 이 차이가 이후 판단의 기준이 되었다.
``0.015보다 확실히 크게 변해야 효과가 있다''고 보고, 그보다 작은 변화는 운과
구별하지 않으려 했다.

점수는 생각보다 낮았다. 사진을 봤을 때는 각 지역의 특징이 드러난 사진이 꽤
있었는데, 부산(0.382 / 0.386)과 오사카(0.370 / 0.363)는 도쿄(0.641 / 0.618)와
타이베이(0.559 / 0.521)보다 지역적 특색이 잘 안 드러난 것 같았다. 부산은 사진으로
보았을 때 부산 바다라기보다는 바다 사진의 느낌이 강했고, 오사카는 일본 표지판이나
간판 거리 사진은 있었지만 오사카 사진이라고 판단하기는 어려웠을 것이라고 생각한다.
부산(5,191장)과 오사카(7,200장)의 사진 수가 적다는 점도 이유가 될 수 있다.

\subsection*{street: 길거리 사진을 더 자주}

\textbf{바꾼 것.} \texttt{sample\_weights}가 사진이 나온 곳에 따라 뽑히는 횟수를
정하도록 했다: commons $\times$1, youtube $\times$2, osv5m $\times$4. 사진은 여전히
44,257장 전부이고, 변환·최적화·스케줄은 시작 레시피 그대로다.

\textbf{왜.} baseline의 틀린 사진 대부분이 길거리 장면이었고, 검증 사진은 학습보다
길거리 사진 비중이 높았다(학습 osv5m 7\%, 검증 14\%). 모델이 랜드마크 말고 평범한
거리에서 단서를 찾게 하고 싶었다. 부산에는 osv5m 사진이 0장이라 오히려 부산이 더
안 맞을 수도 있다고 예상했다.

\textbf{결과.} 0.4857과 0.4812, 평균 0.4835로 baseline 평균 0.4954보다 0.012
낮다. 여전히 0.015 안이라 확실히 나빠졌다고도 말할 수 없다.

\textbf{내가 생각한 의미.} 점수가 오르지는 않고 오히려 조금 떨어졌다. 도시별로는
도쿄($-0.022$)와 타이베이($-0.023$)가 가장 많이 떨어졌다. 길거리 사진을 늘린 만큼
commons 사진을 보는 비율이 64\%에서 43\%로 줄었는데, 도쿄는 도쿄타워·스카이트리,
타이베이는 중정기념당·사원처럼 commons에 있는 랜드마크 사진이 도시를 맞히는 데
도움이 됐던 것 같다. 또 검증 사진도 66\%가 commons라서, 길거리 사진을 늘린 효과가
검증 점수에는 잘 드러나지 않았을 수 있다. 그리고 osv5m 사진은 3,110장밖에 없는데
4배 자주 보여줬기 때문에, 새로운 단서를 배우기보다 같은 사진을 반복해서 외웠을
수도 있다. 부산은 시드 0에서는 떨어지고 시드 1에서는 오히려 가장 높게 나와서
흔들림이 커서 판단하기 어려웠다. 결론적으로 사진을 ``얼마나 자주 보여주는지''
바꾸는 것만으로는 부족하고, 다음에는 길거리와 관계없는 사진을 아예 빼보기로 했다.

\subsection*{drop\_street: 길거리와 상관없는 사진 빼기}

\textbf{바꾼 것.} \texttt{recipe.py}를 시작 레시피로 되돌려(가중치 없음,
\texttt{sample\_weights}가 \texttt{None}) recipe\_sha가 baseline과 같은
\texttt{17ac485223fa}가 되게 하고, 학습 사진만 44,257장에서 37,293장으로 줄였다.
그래서 baseline과의 차이는 오직 뺀 사진이다. 어떤 사진을 뺐는지는 다음 절에 있다.

\textbf{왜.} 검증 사진과 리더보드 사진, 최종 사진이 모두 평범한 길거리 장면인데
학습 사진에는 유물, 음식, 기자회견, 모터쇼 차처럼 길거리에서 볼 수 없는 사진이 많이
섞여 있었다. 학습 횟수가 정해져 있으니 이런 사진을 빼면 모델이 길거리 사진을 더
많이 보고 길거리에서 도시를 구별하는 단서를 더 잘 배울 것이라고 생각했다. 서울과
타이베이에 유독 유물·박물관 사진이 많아서, 모델이 유물 자체보다 어두운 배경에
조명이 비춰지는 물건이라고만 외웠을 수도 있다고 보았다.

\textbf{결과.} 0.4841과 0.4918, 평균 0.4880으로 baseline보다 0.007 낮다.

\textbf{내가 생각한 의미.} 전체 점수는 0.495에서 0.488로 떨어졌는데 0.015보다
작아서 전체적으로 효과가 있었다고 말할 수 없다. 하지만 도시별로는 두 시드가 같은
방향으로 움직였다. 부산과 도쿄는 여섯 번의 실험 중 가장 높았고, 오사카와 서울은
가장 낮았다. 부산은 모터쇼 자동차, 화재 사건, 인물 사진처럼 부산과 상관없는 사진이
빠지면서 실제 부산 길거리의 특징을 더 잘 배운 것 같다. 반대로 서울과 타이베이는
예상과 달리 떨어졌는데, 이 두 도시는 사진이 가장 많이 빠진 도시였다(서울 15.7\%,
타이베이 26.9\%). 유물이나 인물 사진에도 한글, 한자, 한복처럼 나라를 알려주는 단서가
있었는데 그걸 같이 잃었을 수 있다. 오사카가 떨어지고 도쿄가 오른 것은, 오사카
사진이 18.7\% 빠지면서 오사카가 더 적은 도시가 되었고, 모델이 헷갈리는 사진을 도쿄
쪽으로 더 많이 보낸 것이라고 생각한다. 원래 가장 큰 실수가 오사카를 도쿄로 부르는
것이었기 때문이다. 사진을 빼는 것은 도시마다 효과가 달라서, 전체 평균만 보면 효과가
서로 상쇄되어 보이지 않는다.

\subsection*{balanced: 다섯 도시를 똑같이}

\textbf{바꾼 것.} \texttt{sample\_weights}에서 도시마다 똑같은 몫을 주도록 했다:
사진 하나의 무게가 그 도시 사진 수의 역수이므로, 부산의 5,191장이
서울의 11,274장과 같은 횟수만큼 뽑힌다. 사진은 44,257장 전부이고 레시피의 나머지는
시작 레시피 그대로다.

\textbf{왜.} 두 번의 실험에서 같은 결론이 나왔다. drop\_street에서 오사카 사진이
줄자 오사카 점수가 떨어지고 도쿄 쪽으로 쏠렸고, baseline의 부산·오사카는 사진 수가
적은 도시라 계속 낮았다. 그래서 다음에는 사진이 적은 부산과 오사카를 더 자주
보여주는 방법을 시도해보고 싶었다.

\textbf{결과.} 0.5038과 0.5123, 평균 0.5081로 baseline 평균 0.4954보다 0.013
높다. 여덟 번의 실행 중 두 시드가 모두 baseline 두 시드보다 높게 나온 첫 실험이다.
다만 0.013은 아직 0.015보다 작아서 ``확실하다''고 말하기는 어렵다.

\textbf{내가 생각한 의미.} 도시별로는 변화가 아주 컸다. 부산($+0.075$)과
오사카($+0.084$)가 크게 올랐고, 도쿄는 0.104나 떨어졌다. 사진이 적은 부산과
오사카를 더 자주 보여주자, 모델이 헷갈리는 사진을 도쿄나 서울로 보내던 습관이
줄어든 것 같다. drop\_street에서 오사카가 줄자 도쿄로 쏠렸던 것과 반대 방향의
결과라서, 도시별 사진 수가 모델의 판단에 큰 영향을 준다는 걸 알 수 있었다.

점수는 다섯 도시 정답률의 평균이라, 잘하던 도쿄를 조금 포기하고 못하던
부산·오사카를 끌어올리는 것이 전체 점수에 유리했다. 다만 도쿄는 가장 잘 맞히던
도시에서 두 번째나 세 번째로 잘 맞히는 도시가 되었고, 서울과 거의 비슷해졌다.
그래서 완전히 똑같이 맞추기보다 중간 정도로 균형을 맞추는 게 더 좋을 수도 있다.

\begin{center}\footnotesize
\begin{tabular}{lrrrr}
\toprule
도시 & baseline & street & drop\_street & balanced \\
\midrule
seoul & 0.5569 & 0.5592 & 0.5418 & 0.5217 \\
busan & 0.3838 & 0.3773 & 0.4011 & 0.4587 \\
tokyo & 0.6295 & 0.6076 & 0.6456 & 0.5258 \\
osaka & 0.3666 & 0.3562 & 0.3365 & 0.4506 \\
taipei & 0.5399 & 0.5170 & 0.5147 & 0.5834 \\
\midrule
전체 & 0.4954 & 0.4835 & 0.4880 & 0.5081 \\
\bottomrule
\end{tabular}

{\footnotesize 표 2. 도시별 검증 점수, 두 시드 평균. 전체는 표 1의 평균과 같다.}
\end{center}

\subsection*{최종 레시피와 선택}

최종 Track A 레시피는 \texttt{balanced}다. 바뀐 것은 \texttt{sample\_weights}
하나뿐이고, 선택은 학습 사진 44,257장 전부다(\texttt{selection.csv}가 없다). 이
레시피는 release \texttt{track-a-3}, ``every city drawn equally often''으로 냈다.
검증 점수는 0.5038(시드 0)과 0.5123(시드 1)이고, 리더보드 점수는
\fillin{리더보드 점수와, 검증 점수와 다르다면 그 이유}.

% ---------------------------------------------------------------------------
\section*{Decisions you made from photos}

\begin{figure}[h]
\centering
\picturefile[0.92\linewidth]{figures/mistakes_baseline_seed0_busan_00.png}
\caption{baseline(시드 0)이 부산 사진이라고 확신했지만 틀린 48장. 검증 사진을 두 장에
한 장씩 채점해 도시마다 가장 확신한 오답을 모은 것이다.}
\end{figure}

\textbf{무엇을 봤나.} \texttt{curate.py mistakes}로 만든 시트 5장(도시마다 1장)을 다
봤다. 검증 사진 6,039장을 채점했고, 도시별 오답은 서울 549/1218, 부산 434/695,
도쿄 677/1940, 오사카 677/1064, 타이베이 494/1122였다.

\textbf{공통점 두 가지.} 첫째, 틀린 사진 대부분이 같은 나라의 다른 도시로 갔다.
오사카 사진은 거의 다 도쿄로, 부산 사진은 대부분 서울로 갔다. 모델이 한국인지
일본인지는 어느 정도 구별하지만 같은 나라 안에서 도시를 나누는 단서는 잡지 못한 것
같다. 둘째, 학습 사진을 봤을 때는 랜드마크와 관광지가 많았는데, 틀린 사진을 모아
놓으니 차 안에서 찍은 도로, 영상 캡처, 360도 사진처럼 평범한 길거리 장면이
대부분이었다. 모델은 랜드마크로 도시를 구별하는 법을 배웠는데 검증 사진에는 그런
단서가 거의 없어서 틀린 것 같다. 이 관찰에서 Track A의 두 실험(street,
drop\_street)이 나왔다.

\textbf{행동하기 전에 적은 규칙.} ``길거리에서 볼 수 없는 장면은 뺀다.''
\begin{enumerate}
\item 박물관·전시물: 유물, 도자기, 조각, 설명 패널, 전시장 내부
\item 음식·음료 클로즈업: 접시, 음식 사진, 술병
\item 실내 행사: 기자회견, 회의, 시상식, 강연, 공연장, 단체 기념사진
\item 행사장·전시장 자동차: 모터쇼 차, 차 실내, 바퀴 클로즈업
\item 동식물 클로즈업: 새, 꽃, 고양이·강아지, 동물원·수족관
\item 인물 사진: 코스프레, 모델 사진처럼 사람이 화면 대부분인 것
\item 글자만 있는 사진: 문서, 표지판 클로즈업 없이 종이·포스터만 있는 것
\item 특정 사건·행사 하나를 수십 장씩 찍은 묶음 (화재 현장, 의전 행사 등)
\end{enumerate}
사진을 고르는 동안 이 규칙에 세 가지를 더했다: 돌 기념물과 명판의 클로즈업, 한
사건이나 행사를 찍은 묶음(화재, 훈련, 의전 행사, 축구 경기, 불꽃놀이), 항공기와
공항. 반대로 길거리 풍경, 역 표지판, 건물 외관, 블랙박스와 360도 사진은 남겼다.

\textbf{44,257장에서 어떻게 찾았나.} 사진을 한 장씩 볼 수는 없어서, 같은 곳에서
나온 사진을 묶음으로 봤다. YouTube 영상 하나의 프레임과 Commons 업로더 한 사람의
사진은 한 장면 안에서 찍힌 것이라, 묶음 하나를 판단하면 사진 수십 장이 한꺼번에
정해진다. 114개 영상(12,900프레임)과, 사진이 20장 이상인 Commons 업로더 284명
(19,811장)으로 398장의 시트를 만들었고, 큰 묶음부터 봤다. 이것으로 학습 사진
44,257장 가운데 32,711장(74\%)을 덮었다.

\begin{figure}[h]
\centering
\picturefile[0.92\linewidth]{figures/grp_busan_commons_05_00.png}
\caption{묶음 하나가 시트 한 장. 부산의 업로더 한 사람이 올린 사진이고, 화재 현장
한 사건을 찍은 묶음이라 규칙 8번으로 통째로 뺐다. 큰 묶음부터 보도록 정렬되어 있다.}
\end{figure}

\outputfile{outputs/grouped_sheets.txt}

\outputfile{outputs/bundle_totals.txt}

\textbf{무엇이 바뀌었나.} 규칙에 해당하는 사진 6,964장을 뺐고, 학습 사진은 44,257장
에서 37,293장이 되었다. 묶음 단위로 뺀 것이 94묶음(6,798장), 시트 번호로 고른 것이
9시트(166장)이고, 같은 사진이 두 번 지목된 적은 없다. \texttt{curate.py count}의
결과는 다음과 같다.

\outputfile{outputs/curate_count_drop_street.txt}

\textbf{점수 전과 후.} 뺀 사진만 바꾸고 레시피는 baseline으로 되돌린
drop\_street는 0.4841(시드 0)과 0.4918(시드 1)로, baseline의 0.5028과 0.4879에서
평균 0.007 내려갔다. 전체 점수는 움직이지 않았지만 도시별로는 부산과 도쿄가 올랐고
서울·오사카·타이베이가 내려갔다(표 2).

\textbf{에이전트에게 보낸 요청.} 아래는 그 결정을 위해 내가 그대로 보낸 요청이다.

\begin{Verbatim}[fontsize=\footnotesize, breaklines, breakanywhere, frame=single, framesep=2mm]
Drop scenes you can't see on the street: museum objects and panels, food
and drink close-ups, indoor events (press conferences, meetings,
ceremonies, performances), cars at motor shows, close-ups of animals and
flowers, portraits where a person fills the frame, and documents or
posters only. Keep street views, station signs, building exteriors,
dashcam and 360 photos.
\end{Verbatim}

\begin{Verbatim}[fontsize=\footnotesize, breaklines, breakanywhere, frame=single, framesep=2mm]
My drop rule now also covers: close-ups of stone monuments and plaques;
bundles of one event or incident (fire, drills, ceremonies, football
matches, fireworks); aircraft and airports.

For the next experiment, put recipe.py back to the starter recipe (no
street weighting), so the only change from baseline is the dropped
photos. Tag it "drop_street".
\end{Verbatim}

\begin{Verbatim}[fontsize=\footnotesize, breaklines, breakanywhere, frame=single, framesep=2mm]
Drop the WHOLE bundle (use .paths.txt) for these:
osaka_commons_01 (food), 02 (animals), 09 (food), 11 (monument plaques),
12 (motor show cars), 21 (flowers), 25 (flowers, animals), 34 (flowers,
animals, food), 40 (cars, animals)
osaka_youtube_03 (selfie faces), 11 (cars), 12 (shoes, food), 13 (figures,
food), 15 (selfie faces)
busan_commons_01 (food), 02 (portraits), 03 (cars), 05 (one event: fire),
12 (portraits), 13 (one event, portraits), 24 (museum objects),
26 (portraits), 31 (animals), 32 (museum objects, paintings),
33 (museum objects)
seoul_commons_03 (cars), 05 (museum), 06 (food), 13, 15, 17 (portraits),
18 (museum), 21 (one event: rescue drill), 23 (food), 26 (indoor talks),
33 (museum), 34 (portraits), 39 (one event: palace ceremony), 42,
43 (museum), 44, 53 (portraits), 47, 62 (aircraft), 59 (portraits,
indoor), 63 (animals), 65 (food), 67 (portraits), 69 (museum, medals),
72 (animals), 73 (flowers), 75 (trees), 76 (objects)
taipei_commons_01 (portraits, indoor), 04 (food), 06, 07 (museum objects),
09 (portraits), 11 (food, indoor), 19, 21, 23 (portraits), 20, 26 (museum
objects), 22 (portraits, characters, food), 29, 31, 34, 35, 38, 49, 54,
55 (portraits), 33 (portraits, cars), 42 (food, indoor), 44 (portraits,
animals), 51 (animals, plants)
tokyo_commons_06 (fireworks), 09 (food), 17 (fireworks, animals,
portraits), 18, 25, 59 (portraits), 20 (cars), 29, 30, 37, 65, 66
(aircraft), 40, 42 (animals, plants), 48 (food, fireworks), 49 (museum,
indoor), 64 (food)
\end{Verbatim}

\begin{Verbatim}[fontsize=\footnotesize, breaklines, breakanywhere, frame=single, framesep=2mm]
Drop ONLY these numbers:
osaka_commons_36: 19-26, 29-34, 36 (food, objects, flowers)
osaka_commons_37: 9-21, 31, 35 (aquarium, flowers)
osaka_commons_47: 7-18 (aquarium, indoor restaurant)
busan_commons_16: 9, 11-14, 16, 18-32, 34, 35, 44, 45 (indoor altars,
museum objects, monument close-ups)
busan_commons_22: 5, 6, 12, 14-17, 19-24, 31 (indoor temple halls, plaque
and statue close-ups)
busan_commons_27: 1, 2, 5-13, 15-20, 23 (grave and monument plaques,
portrait)
seoul_commons_20: 1-6, 12-25, 27, 30, 32, 36, 37, 44-48 (indoor meetings,
museum figures)
seoul_commons_36: 1, 2, 19, 20, 28, 32, 33, 35-43 (food and flower
close-ups)
tokyo_commons_07: 28-48 (one event: football match)
\end{Verbatim}

묶음 이름은 시트 이름(\texttt{grp\_<도시>\_<출처>\_<번호>})에서 도시와 출처를 줄인
것이고, 묶음 전체를 뺄 때는 그 시트의 \texttt{.paths.txt}에 있는 사진을 전부 뺐다.
시트 번호로 고른 것은 그 시트에 그려진 사진 중에서 고른 번호만 뺐다. 같은 사진이 두
목록에 겹쳐 지목된 경우는 없었다.

% ---------------------------------------------------------------------------
\section*{Track B}

\fillin{What you tried, one step at a time, and the validation score after
each step.}

\fillin{How you collected any extra photos.}

\fillin{Your validation score next to your leaderboard score, and why they
might disagree.}

\subsection*{Your submission, as the grader sees it}
\outputfile{outputs/check_submission.txt}

% ---------------------------------------------------------------------------
\section*{What made the difference}

\textbf{Track A에서 가장 도움이 된 것.} \texttt{balanced}다. 사진을 한 장도 바꾸지
않고 \texttt{sample\_weights}만 바꿨는데 평균이 0.4954에서 0.5081로 0.013 올랐고,
여덟 번의 실행 중 두 시드가 모두 baseline 두 시드보다 높게 나온 유일한 실험이다.
도시별로는 부산이 0.3838에서 0.4587로, 오사카가 0.3666에서 0.4506으로 올랐다. 두
도시 모두 사진 수가 가장 적은 도시라는 점이 원인이었다고 본다.

\textbf{가장 도움이 안 된 것.} \texttt{street}다. 평균이 0.4835로 0.012 내려갔고
도쿄와 타이베이가 크게 떨어졌다. 사진을 ``얼마나 자주 보여주는지'' 바꾸는 것만으로는
부족하다는 결론이었다. \texttt{drop\_street}는 평균으로는 0.007 내려갔지만, 뺀
사진이 가장 많았던 서울과 타이베이가 내려간 반면 부산과 도쿄는 올라서, 같은 규칙이
도시마다 반대로 작용했다. 그래서 전체 평균 하나로는 판단할 수 없는 변화였다.

\textbf{Track B.} \fillin{Track B에서 가장 도움이 된 변화와 가장 안 된 변화, 그 숫자.}

\textbf{다음 주가 있다면.} 먼저 \texttt{balanced}의 정도를 조절해보고 싶다. 지금은
다섯 도시의 몫을 완전히 같게 했는데, 그러면 가장 잘 맞히던 도쿄의 점수를 0.104나
내주게 된다. 도시 몫을 부분적으로만 균등하게 하는 중간 지점이 더 좋을 수 있다.
다음으로 \texttt{drop\_street}와 \texttt{balanced}를 함께 써보고 싶다. 두 실험은
부산·오사카를 올린 방향이 같으므로 같이 쓰면 더 오를 수 있고, 대신 서울·타이베이의
손실을 얼마나 줄일 수 있는지가 관건이다. Track B에서는
\fillin{Track B에서 다음 주에 할 일}.

% ---------------------------------------------------------------------------
\section*{All runs}
\outputfile{outputs/experiments.csv}

\section*{Lab notebook}
\outputfile{outputs/LOG.md}

\end{document}
